
\subsection*{Abstract}
Real-world Knowledge Graph deployments must ingest continuously arriving data streams --- sensor readings, financial ticks, event logs, and scientific observations --- and immediately make this data available for reasoning. We present CEREBRUM's \textbf{Streaming Engine}, a composable pipeline that transforms heterogeneous continuous signals into typed graph edges through a family of stateless discretizers, then integrates the resulting edges into the live reasoning graph with zero downtime. The pipeline supports five discretizer types (threshold, STDP, delta, windowed-frequency, and pattern) and composes with the existing \texttt{StreamAdapter} and \texttt{IngestionPipeline} layers. In v2.71.0 (Phase 53), the engine gains \textbf{adaptive search strategy}: local graph density is measured per-query and dynamically adjusts beam width, depth limit, and branching factor. In Phase 54, the \texttt{/build} endpoint enables hot CSV reload without server restart. CORS middleware and request-timing middleware are added for production observability. The streaming engine now operates at the same production maturity level as the batch reasoning stack.

\subsection*{1. Introduction}
Traditional Knowledge Graph pipelines assume a static or batch-updated graph: data arrives in bulk, the graph is rebuilt or updated offline, and queries run against a stable snapshot. This model fails for applications where the latency between observation and reasoning must be sub-second: industrial sensor networks, financial surveillance, genomic streaming sequencers, and real-time intelligence fusion.

CEREBRUM's Streaming Engine bridges this gap by providing a composable pipeline from raw signal to typed graph edge to CSA-weighted reasoning path --- all in a single continuous flow with no offline reconstruction step.

\subsection*{2. Streaming Discretizers}
Discretizers transform continuous input signals into discrete symbolic edges. The core discretizers include:

\begin{itemize}
  \item \textbf{Threshold Discretizer}: Continuous float stream $\to$ Edge emitted when value crosses threshold $\theta$.
  \item \textbf{STDP-Discretizer}: Spike/event timestamps $\to$ Directional \texttt{CAUSES} edge based on temporal co-occurrence ($\Delta t$).
  \item \textbf{Delta-Discretizer}: Rate-of-change signal $\to$ Edge emitted when $|\Delta x / \Delta t| \geq \theta_{rate}$.
  \item \textbf{Windowed-Frequency-Discretizer}: Event counts per window $\to$ Edge emitted when co-occurrence frequency exceeds $f_{min}$.
  \item \textbf{Pattern-Discretizer}: Symbolic event sequence $\to$ Edge emitted when pattern match probability $\geq p_{match}$.
\end{itemize}


Each discretizer maintains a small internal sliding buffer and is stateless with respect to the adapter graph. This isolation ensures discretizer failures cannot corrupt the persistent Knowledge Graph.

\subsection*{3. StreamAdapter and IngestionPipeline Integration}
Discretized edges flow into the \texttt{StreamAdapter}, which exposes the \texttt{GraphAdapter} interface and applies the standard \texttt{IngestionPipeline} normalization stack (entity dedup, relation normalization, confidence/provenance assignment) before materializing edges into the live graph. Community membership for new nodes is initialized via a lightweight single-node DSCF assignment (attaching to the nearest existing community centroid) without triggering a full global rebalance.

\subsection*{4. STDP Spike Processing and Causal Significance Filtering}
The \texttt{STDPDiscretizer} implements Hebbian-inspired causal edge inference: when two event streams exhibit consistent temporal co-occurrence within a configurable window $\Delta t_{max}$, a directional \texttt{CAUSES} edge is materialized. Two structural holes in this process were identified and patched:

\begin{itemize}
  \item \textbf{\texttt{min\_causal\_span}}: Blocks materialization when all co-occurrences fall within a burst window shorter than the configured span, preventing adversarial spike floods.
  \item \textbf{\texttt{use\_chi\_squared}}: Applies a chi-squared uniformity test to inter-event intervals; non-uniform (burst) distributions are rejected at $p < 0.05$.
\end{itemize}


These protections are documented in detail in Paper 16 (Production Hardening).

\subsection*{5. Recent Advances (v2.51.1 -> v2.71.0)}

\subsubsection*{5.1 Adaptive Search Strategy (Phase 53)}
The most consequential algorithmic advance in the streaming engine is the introduction of \textbf{adaptive search strategy} (Phase 53). Prior versions used fixed beam parameters (width, depth, branching factor) configured at startup. In streaming graphs, however, local density varies dramatically: a newly-ingested event cluster may produce a dense sub-graph that overwhelms a narrow beam, while a sparsely-observed sensor region starves a wide beam.

Phase 53 adds a \textbf{density-aware parameter selector} that measures local graph density at the traversal entry point before each query:

$$\rho(v, r) = \frac{|\mathcal{E}(B(v,r))|}{|\mathcal{V}(B(v,r))|}$$

where $B(v, r)$ is the $r$-hop ball around the query root $v$. Based on $\rho$, the traversal parameters are selected from a configurable density-to-params map:

\begin{table}[h]
\small\centering
\begin{tabular}{llll}
\toprule
Density Regime & Beam Width & Depth Limit & Branch Factor \\
\midrule
Sparse ($\rho < 1.5$) & 8 & 6 & 4 \\
Normal ($1.5 \leq \rho < 4.0$) & 5 & 4 & 3 \\
Dense ($\rho \geq 4.0$) & 3 & 3 & 2 \\
\bottomrule
\end{tabular}
\end{table}


This dynamic adjustment reduces average query latency by 31\% on streaming graphs with heterogeneous density profiles, while maintaining H@10 within 2\% of the fixed-wide-beam configuration.

\subsubsection*{5.2 Hot CSV Reload via /build Endpoint (Phase 54)}
The \texttt{/build} REST endpoint accepts a multipart CSV upload and triggers a full graph rebuild from the new data --- re-running IngestionPipeline, EmbeddingEngine, StructuralEncoder, and CommunityEngine --- without restarting the server process. During rebuild, in-flight queries continue against the current graph snapshot (via Query Snapshot Isolation). The new graph atomically replaces the old one upon rebuild completion.

This enables continuous graph evolution workflows: operators can upload corrected or enriched CSVs to a running production server and immediately observe the updated reasoning behavior.

\subsubsection*{5.3 Production Middleware (Phase 54)}
Two middleware layers are now applied to all API endpoints:

\begin{itemize}
  \item \textbf{CORS middleware}: Configurable origin allowlist enables the Studio and Dashboard frontends to call the API from browser contexts without proxy configuration.
  \item \textbf{Request-timing middleware}: All requests are annotated with \texttt{X-Process-Time} response headers and structured log entries, enabling latency monitoring via the \texttt{/logs} ring buffer and external APM tools.
\end{itemize}


Together with the \texttt{/logs} endpoint and dashboard.html (Paper 12), these middleware layers provide end-to-end production observability without requiring external infrastructure.

\subsection*{6. Conclusion}
The CEREBRUM Streaming Engine in v2.71.0 has matured from a laboratory prototype into a production-grade continuous ingestion and reasoning pipeline. The adaptive search strategy (Phase 53) brings intelligent runtime adaptation to beam traversal, the \texttt{/build} endpoint (Phase 54) enables zero-downtime graph evolution, and the production middleware stack provides the observability required for enterprise deployment. Combined with the 2,261-test suite and CORS/timing instrumentation, the streaming engine is now a first-class production component of the CEREBRUM framework.